% Complete substitution and finite-closure proofs for Kamp III.2, Lemma 3.

\section{Kamp's Boolean and proper-substitution closure}
\label{sec:annex-p09-kamp-iii2-lemma3}

\begin{annexitem}{annex:p9.kamp.proper.substitution}{thm:p09-kamp-proper-substitution}{Capture avoidance, substitution semantics and the D-translation}
The proper substitution is that of Kamp's Definition~1 on printed p.~40;
the expressibility closure R4 is stated across printed pp.~59--60
\cite[ch.~III, pp.~40--42, 59--60]{Kamp1968}. We prove the needed claim
with its complete first-order and temporal inductions.

\begin{proof}
Fix a nonempty domain \(T\), a binary relation \(<\) on \(T\), and
interpretations of all unary predicate letters. Denote the resulting
structure by \(\mathcal M\). An assignment \(g\) maps the variables
of \(L\) into \(T\). The notation \(g[x\mapsto a]\) changes its value
at \(x\) to \(a\) and leaves every other value unchanged.

\paragraph{Variable dependence and bound renaming.}
Truth of an \(L\)-formula depends only on the assignment to its free
variables. For a unary atom this is its argument; for an order atom it
is its two arguments. Negation preserves this dependence, and conjunction
or disjunction takes the union of the two sets of relevant variables.
For \(\exists x\,\alpha\), assignments agreeing on the free variables
other than \(x\) extend with any same value \(a\) at \(x\); induction
then gives the same truth for \(\alpha\), for every \(a\in T\).
Existence of a successful value agrees. The argument for \(\forall x\)
uses all such values and yields the same conclusion. This proves the
assertion by structural induction, including vacuous quantifiers.

If a quantifier's bound variable \(x\) is renamed to a fresh \(z\),
with exactly the occurrences bound by that quantifier renamed, giving
\(z\) the old value of \(x\) preserves the values of all affected
atoms. The same structural argument preserves the matrix's truth. The
outer quantifier still ranges over the same set \(T\), so its truth is
unchanged for either quantifier. Successive such renamings prove semantic
invariance under alphabetic variation. With no capture of a free argument,
the same argument also gives
\[
 \mathcal M,g\models\theta[y/v]
 \quad\Longleftrightarrow\quad
 \mathcal M,g[v\mapsto g(y)]\models\theta
\tag{1}
\]
for any formula \(\theta\), with its bound variables chosen fresh
from \(y\). We will use it in particular when \(v\) is the sole
free variable. The case \(y=v\) is the identity. In (1), the renamed
occurrences are free occurrences;
bound occurrences of the same spelling are first separated by alphabetic
variation.

To verify (1) structurally, an argument occurrence of \(v\) on the
right has value \(g(y)\), exactly the value of its renamed argument on
the left. Every unchanged atomic argument has the same value on both
sides. This proves the unary and order-atom cases. Negation, conjunction
and disjunction preserve the resulting equivalences. At a quantifier
\(\exists z\) or \(\forall z\), choose \(z\) different from
\(v,y\) by the allowed bound renaming. Updating the assignment at
\(z\) then commutes with updating it at \(v\), and it does not
change \(g(y)\). The matrix induction equivalence holds for every
assigned value of \(z\). Taking existence or universality proves the
two quantifier cases of (1).

\paragraph{Why fresh copies exist and agree with the source operation.}
Each formula contains finitely many variable names, whereas the list
\(t_0,t_1,\ldots\) is infinite. Rename bound variables successively,
using different names outside the finite set of all variables of
\(\phi\), the free parameter of \(\psi\), and the names already
chosen. Thus a suitable alphabetic variant exists. It occurs somewhere
in the source's enumeration, so the first eligible variant exists too.

Any two eligible variants can have their bound names renamed further to
one common fresh list. After insertion at an original \(P_i(y)\),
their free argument occurrences are both precisely \(y\); the same
common renaming identifies their bound occurrences. If a source-selected
variant used its free parameter's spelling also at a bound occurrence,
separate those bound occurrences first. This changes only their bound
names and does not change the inserted free argument. Doing this in every
copy shows that any direct fresh-copy implementation gives an alphabetic
variant of the result selected by the source convention. The first entry
of the enumeration is not needed by the temporal construction.

\paragraph{First-order substitution and a modified unary interpretation.}
Let \(\psi\) have exactly the free variable \(v\). Define
\[
 A\defeq\{a\in T:\mathcal M,g[v\mapsto a]\models\psi\}.
\]
The free-variable dependence just proved shows that this set is independent
of the choice of ambient assignment \(g\). Let \(\mathcal M_i^A\)
have the same domain, the same binary relation, and the same unary
interpretations as \(\mathcal M\), except that \(P_i\) is interpreted
by \(A\). In particular \(\psi\) is evaluated in the original
\(\mathcal M\) to define \(A\), even when \(P_i\) occurs in
\(\psi\). This is a definition of one set, with no recursive equation.

For an arbitrary \(L\)-formula \(\alpha\), we claim
\[
 \mathcal M,g\models\operatorname{psub}_{P_i}(\alpha,\psi)
 \quad\Longleftrightarrow\quad
                         \mathcal M_i^A,g\models\alpha.
\tag{2}
\]
Use fresh copies of \(\psi\) whose bound names avoid every variable
of \(\alpha\); the preceding paragraph transfers the result to the
source-selected copies. Prove (2) structurally on the original formula
\(\alpha\), without recursively visiting an inserted copy.

For \(\alpha=P_i(y)\), equation (1) says that the inserted formula is
true exactly when \(g(y)\in A\), which is the interpretation of this
atom in \(\mathcal M_i^A\). For \(\alpha=P_j(y)\) with \(j\ne i\),
the atom is unchanged and the two structures give \(P_j\) the same
interpretation. For an order atom \(x<y\), both the atom and the binary
relation are unchanged. These exhaust the atoms of the equality-free
language.

For \(\neg\alpha\), the induction equivalence negates on both sides.
For \(\alpha\land\beta\), truth is the conjunction of the two child
truths, so their two induction equivalences give (2). For
\(\alpha\lor\beta\), truth is the disjunction of the two child
truths, and the same two equivalences prove (2) with disjunction. Other
Boolean abbreviations may first be expanded into these connectives.

For \(\exists x\,\alpha\), choose inserted bound names fresh from
the entire formula, including \(x\). The resulting outer quantifier
still binds the argument occurrences corresponding to the original atoms
whose argument was \(x\). Under each assignment \(g[x\mapsto a]\),
the induction hypothesis gives the equivalence of the two matrices.
The set \(A\) is the same for every such assignment, because \(\psi\)
has no additional free parameter. Taking existence over \(a\in T\)
proves (2). For \(\forall x\,\alpha\), apply the same equivalence
at every \(g[x\mapsto a]\) and take universal quantification over
\(a\in T\). This proves the universal case and completes (2).

The same induction checks free variables syntactically. Each replaced
atom \(P_i(y)\) becomes a formula with exactly \(y\) free, since
\(\psi\) has exactly \(v\) free. The other atoms are unchanged.
Boolean nodes take the same union of free variables as before, and a
quantifier removes exactly its bound argument variable. Fresh inserted
bound names remove no other variable. Thus proper substitution preserves
the complete free-variable set of its outer formula; if that set is
\(\{w\}\), it remains exactly \(\{w\}\).

\paragraph{Temporal substitution and the corresponding interpretation.}
For a temporal formula \(\zeta\), put
\[
                 B\defeq\{a\in T:\mathcal M,a\models\zeta\}.
\]
For every temporal formula \(\eta\) and \(t\in T\), we prove
\[
 \mathcal M,t\models\eta[\zeta/q_i]
 \quad\Longleftrightarrow\quad
                           \mathcal M_i^B,t\models\eta.
\tag{3}
\]
At the atom \(q_i\), the left side is the truth of the supplied
\(\zeta\) and hence membership of \(t\) in \(B\), exactly the
right side. At \(q_j\) for \(j\ne i\), the proposition's unary
interpretation is unchanged. Negation reverses the equivalent child
truths, and conjunction takes their conjunction. Disjunction, when used,
is the declared abbreviation and therefore has the same property.

For \(S(\alpha,\beta)\), its substituted version is true at \(t\)
exactly when some \(s<t\) satisfies \(\alpha[\zeta/q_i]\) in
\(\mathcal M\) and every \(r\) with \(s<r<t\) satisfies
\(\beta[\zeta/q_i]\). Apply the induction hypothesis at \(s\)
and separately at each such \(r\). This gives the truth of
\(\alpha\) at the same witness and of \(\beta\) throughout the
same open interval in \(\mathcal M_i^B\), which is precisely the
right side of (3). For \(U(\alpha,\beta)\), the witness instead
satisfies \(t<s\) and the interval points satisfy \(t<r<s\).
The same two applications of induction at these actual future points
give its strict-until clause in \(\mathcal M_i^B\). Neither case
adds a point to an empty interval or creates a missing witness. This
completes (3), including endpoints and empty intervals.

\paragraph{Compatibility with the named-variable D-translation.}
Write \(D(\eta)[w/t_0]\) only after renaming its bound variables
fresh from \(w\), as in the body. Here is the full truth verification
needed for this use of \(D\). For \(q_j\), its translation is
\(P_j(w)\), whose truth under \(g\) is the proposition's truth at
\(g(w)\). The negation and conjunction clauses of \(D\) give exactly
the corresponding operations on the child truth values, so their cases
follow from the two child induction assertions.

For \(S(\alpha,\beta)\), choose fresh \(a,b\) as in the displayed
definition of \(D\). Its named-variable image is
\[
 \exists a\bigl(a<w\land D(\alpha)[a/t_0]
    \land\forall b\,(a<b<w\to D(\beta)[b/t_0])\bigr).
\]
By (1) and the child induction assertions, the first translated child
is true exactly when \(\alpha\) holds at the chosen value of \(a\),
and the second exactly when \(\beta\) holds at the tested value of
\(b\). Thus the formula requires a strict past witness and the same
open-interval truths as \(S(\alpha,\beta)\) at \(g(w)\).
For \(U(\alpha,\beta)\), the translated witness guard is \(w<a\)
and the universal guard is \(w<b<a\). The same child assertions give
truth of \(\alpha\) at that future witness and of \(\beta\) at
each of those interval points, exactly the strict-until semantics.
These are every primitive temporal constructor, so structural induction
identifies truth of \(D(\eta)[w/t_0]\) under \(g\) with temporal
truth at \(g(w)\).

Apply (2) with outer formula
\(D(\eta)[w/t_0]\) and replacement \(D(\zeta)[v/t_0]\).
The set \(A\) in (2) is then exactly \(B\) in (3). Equations
(2)--(3) and the same truth property of \(D\) yield
\[
\begin{gathered}
 \mathcal M,g\models
  \operatorname{psub}_{P_i}
       \bigl(D(\eta)[w/t_0],D(\zeta)[v/t_0]\bigr)\\
 \Longleftrightarrow\\
 \mathcal M,g\models D(\eta[\zeta/q_i])[w/t_0].
\end{gathered}
\tag{4}
\]
This argument holds for every binary relation \(<\) on the nonempty
domain: the two interpretations change only one unary predicate. It proves
semantic compatibility, without claiming that independently selected fresh
names make the two displayed strings literally identical.

Now impose the lemma's linear-order hypothesis and suppose \(\eta\)
expresses \(\phi(w)\) and \(\zeta\) expresses \(\psi(v)\).
For the latter pair, their truth equivalence at every point makes
\(A=B\). The modified structure \(\mathcal M_i^A\) is still a
linear structure with an allowed unary interpretation. The first
expressibility assumption can therefore be applied in that very
structure. For every assignment \(g\),
\[
\begin{aligned}
 \mathcal M,g\models\operatorname{psub}_{P_i}(\phi,\psi)
 &\Longleftrightarrow\mathcal M_i^A,g\models\phi\\
 &\Longleftrightarrow\mathcal M_i^A,g(w)\models\eta\\
 &\Longleftrightarrow\mathcal M,g(w)\models\eta[\zeta/q_i]\\
 &\Longleftrightarrow
       \mathcal M,g\models D(\eta[\zeta/q_i])[w/t_0].
\end{aligned}
\tag{5}
\]
This proves the asserted expressibility and its exact variable condition.

\paragraph{Absent targets, shared letters and finite execution.}
If the target does not occur in the outer first-order formula, its proper
substitution is the identity. Equation (5) still holds even if a supplied
temporal expression contains a redundant occurrence of \(q_i\), because
it expresses a condition independent of that interpretation. In the
class construction below we will use expressions with exactly the
corresponding predicate letters, so an absent first-order target is absent
from its temporal expression too.

If the replacement contains the target, equations (2)--(3) evaluate that
replacement in the original structure, not in a recursively modified
structure. Its occurrence is therefore legitimate. Other shared predicate
letters likewise retain their original interpretations. The temporal
recursion visits the finite original syntax tree and copies the finite
replacement at selected leaves, so it terminates. Repeating finitely many
such steps remains finite; zero steps leave the original formula. The
first-order fresh-name choices and all copies are finite as well.
\end{proof}
\end{annexitem}

\begin{annexitem}{annex:p9.kamp.iii2.lemma3}{thm:p09-kamp-iii2-lemma3}{Every finite E-construction has a temporal expression}
This supplies the Boolean and proper-substitution closure argument used in
Kamp's Lemma~3 \cite[ch.~III, pp.~67, 73]{Kamp1968}.

\begin{proof}
\textbf{The staged grammar is exactly the stated least closure.}\enspace
The union of the stages contains every generator because it contains
stage zero. If two members occur at stages \(a,b\), both occur at
stage \(\max(a,b)\) because stages retain preceding members. Their
conjunction, disjunction, or proper substitution therefore occurs at the
next stage. The negation of a member occurs at the next stage after its
own stage. Thus the union is closed under all four operations.

Conversely, any class containing the generators and closed under these
operations contains stage zero. If it contains stage \(m\), it contains
every retained member and every new constructor at stage \(m+1\).
Induction gives every stage and their union. This proves leastness.
Each membership has a finite construction history, by following the
finitely many children of a constructor back to earlier stages. No
infinitary closure or membership decision claim is needed.

\paragraph{All five kinds of generators.}
Fix the variable \(w\). For an \(\operatorname{At}(w)\) generator,
each positive literal \(P_i(w)\) is expressed by \(q_i\) and each
negative literal \(\neg P_i(w)\) by \(\neg q_i\), directly from
the atomic and negation clauses of \(D\). A nonempty finite conjunction
is expressed by the corresponding conjunction of those literal
expressions: induction on its number of factors gives the same truth
for every factor and hence for the conjunction. In the particular
tautology case \(P_0(w)\lor\neg P_0(w)\), the prescribed
\(\neg(q_0\land\neg q_0)\) is true at every point and expresses
that formula. It uses the same predicate letter.

For a generator in \(\operatorname{Bp}(w)\), Lemma~1 supplies its
actual past recurrence \(L_n\). For a generator in
\(\operatorname{Bf}(w)\), the same lemma supplies its actual future
recurrence \(R_0\). Those proofs include every repeated-letter word,
strict endpoint guard, and empty intervening interval.
For a generator in \(\operatorname{Lp}(w)\), Lemma~2 supplies
\(\Lambda^-_\sigma\), and for a generator in
\(\operatorname{Lf}(w)\), it supplies \(\Lambda^+_\sigma\).
Their stage-zero one-letter cases and every positive finite pattern are
included in that lemma. They use the witness-point limiting definitions,
not the separate handbook current-point comparison. These are all five
generator classes of \(\mathcal G(w)\).

Every generator has exactly \(w\) free, by its stated grammar. Its
chosen temporal expression has exactly the proposition letters
corresponding to its unary predicate letters: this is immediate for
the literal and tautology cases, and follows by inspection of the
already supplied basic and limiting recurrences for the other four.
Negation, conjunction, and the S/U recurrences do not introduce another
letter; every position of each input word occurs in its expression.

\paragraph{Every Boolean step.}
Assume the child expressions have been constructed and verified. If
\(\eta\) expresses \(\phi\), the negation clause of \(D\) and
classical negation show that \(\neg\eta\) expresses \(\neg\phi\).
If \(\eta,\zeta\) express \(\phi,\psi\), respectively, their
truth equivalences hold at the same current point in every allowed
structure. Taking conjunction gives the required equivalence for
\(\eta\land\zeta\) and \(\phi\land\psi\).

For a disjunction step, \(\neg(\neg\eta\land\neg\zeta)\) is
true at a point exactly when at least one of the two child expressions
is true there. By the two child equivalences, this is exactly the truth
of \(\phi\lor\psi\). This proves the disjunction case using only
primitive temporal constructors. Each Boolean result still has exactly
\(w\) free syntactically. Its temporal letter set is the corresponding
unchanged or unioned child letter set. Retention of a previous-stage
member keeps its already verified expression.

The empty Boolean abbreviations in the body require no additional rule.
Their displayed \(T^K_w\) and \(F^K_w\) are source At-generators.
The former was treated above; the latter is expressed by
\(q_0\land\neg q_0\), which is false at every point and uses the
same letter. They retain \(w\) syntactically despite their constant
truth values. No new primitive truth constant or predicate-free temporal
formula is being assumed.

\paragraph{Every proper-substitution step.}
At such a node the two earlier operands \(\phi,\psi\) have exactly
\(w\) free and their expressions \(\eta,\zeta\) have already been
verified. Apply \cref{thm:p09-kamp-proper-substitution} with \(v=w\).
It gives exactly the named-variable equivalence for
\(\operatorname{psub}_{P_i}(\phi,\psi)\) and
\(\eta[\zeta/q_i]\), as well as the retained free-variable set.
The preceding annex item has proved every atom, Boolean, quantifier,
since and until case needed for this application. This invocation is
therefore not an appeal to an unproved substitution principle.

The letter claim is preserved as well. If \(P_i\) is absent from
\(\phi\), the induction claim makes \(q_i\) absent from
\(\eta\), and both substitution operations leave their outer
formulas unchanged. If \(P_i\) occurs, every such original predicate
atom is replaced and every corresponding original proposition occurrence
is replaced. The unary letters of the result are the old letters other
than \(P_i\), together with the letters of \(\psi\); the temporal
letters are the corresponding old letters other than \(q_i\), together
with the letters of \(\zeta\). This remains true if the replacement
contains \(P_i\) or shares another letter with the outer formula:
the union may retain a letter, and there is no second substitution inside
an inserted copy. Fresh renaming affects only individual variables.

Induction on the finite construction therefore proves both the
expressibility equivalence and the letter claim for every member of
\(\operatorname{E}(w)\). The supplied procedure performs finitely
many seed constructions and finitely many structural substitutions, so
its output is a finite \(\mathsf{TL}_1\)-formula. Repeating the
argument for every variable \(t_j\) proves the lemma for the union
\(\operatorname{E}\).

\paragraph{The basic-future example with an inserted limit.}
Choose distinct temporary variables \(u,z,x,r\), all fresh from
\(w\). The outer basic future formula can be written
\[
 \exists u\bigl(w<u\land A(w)
   \land\forall z\,(w<z<u\to B(z))\land C(u)\bigr).
\]
Proper substitution of \(\partial^-C(w)\) for its original \(C\)
occurrence gives, up to harmless bound renaming,
\[
\begin{split}
 \exists u\bigl(w<u\land A(w)
   &\land\forall z\,(w<z<u\to B(z))\\
   &\land\forall x\,(x<u\to
            \exists r\,(x<r<u\land C(r)))\bigr).
\end{split}
\]
The inner universal-existential condition is the past limit at \(u\),
not at \(w\). Lemma~2 expresses it there by \(P'(\widehat C)\).
The strict-until clause consequently identifies the entire display with
\(\widehat A\land U(P'(\widehat C),\widehat B)\) at \(w\),
exactly the temporal substitution prescribed in the body. The predicate
\(C(r)\) inside the inserted limit is not substituted again. If there
is no future endpoint, the existential and the until formula both fail;
if an intervening interval is empty, its universal \(B\)-condition is
vacuous on both sides. The limit factor has the endpoint and adjacency
behavior already proved in Lemma~2. Thus the example imposes no unstated
density or endpoint condition.

The whole proof preserves the underlying linear order and changes only
unary interpretations when proving substitution. It never uses a least
upper bound, greatest lower bound, or any interval-property lemma. The
source's move to \(\mathsf{Com}\) occurs after the lemma just proved;
Lemma~4 and the remainder of Theorem~III.2 are still separate work.
\end{proof}
\end{annexitem}
